% TOP_PROBLEM: {"id": 3000066, "title": "Polyhedral description of kernels", "category_id": 2, "category": {"id": 2, "name": "combinatorics", "display_name": "Combinatorics", "description": "Counting problems, graph theory, discrete structures.", "slug": "combinatorics", "order_index": 2, "created_at": "2026-07-31T15:26:25.671Z"}, "status": "open", "problem_number": "AMR-029-0066", "set_id": 13}
% FIRST_POSTED: 2026-10-02
% ENTRY_KIND: statement_only
% UnsolvedMath ID 3000066; source catalogue code AMR-029-0066
% !TeX program = lualatex
% Definition and sources only; this file is not an LLM solution attempt.
\documentclass[11pt,a4paper]{article}
\usepackage[margin=25mm]{geometry}
\usepackage{fontspec}
\setmainfont{lmroman10-regular.otf}[BoldFont=lmroman10-bold.otf,
 ItalicFont=lmroman10-italic.otf,BoldItalicFont=lmroman10-bolditalic.otf]
\usepackage{amsmath,amssymb,mathtools}
\usepackage{unicode-math}
\setmathfont{latinmodern-math.otf}
\AtBeginDocument{\let\setminus\smallsetminus}
\usepackage{xurl,hyperref}
\hypersetup{colorlinks=true,urlcolor=blue}
\setcounter{secnumdepth}{0}
\setlength{\parindent}{0pt}
\setlength{\parskip}{5pt}
\setlength{\emergencystretch}{3em}
\begin{document}
\section*{Polyhedral description of kernels}\label{3000066}
{\small UnsolvedMath ID 3000066 · AMR-029-0066 · Combinatorics · AMR Open Problem Lists}
\subsection{Definitions and mathematical statement}
For a finite loopless digraph $D=(V,A)$, a kernel is an independent
set $K$ such that every vertex outside $K$ has an outgoing arc into
$K$. Let $\chi^K\in\{0,1\}^V$ be its incidence vector and define
\[
 P_{\rm ker}(D)=\operatorname{conv}\{\chi^K:K\text{ is a kernel of }D\}.
\]
Use the empty polytope when no kernel exists. The underlying undirected
graph joins two vertices whenever there is an arc in either direction.

Determine structural classes of digraphs for which this polytope has
an explicit linear-inequality description. The inequalities must be
described directly from the graph or a specified representation,
rather than obtained by first enumerating all kernels. State whether
the system has polynomially many inequalities, permits polynomial-time
separation, or is only a combinatorial characterization; these are
different strengths of answer.

Every kernel vector obeys
\[
 x\ge0,\qquad x(Q)\le1\quad(Q\text{ a clique}),\qquad
 x_v+\sum_{u:vu\in A}x_u\ge1\quad(v\in V).
\]
Singleton cliques are included. These necessary inequalities suggest
a starting relaxation, but the problem does not assert that they
always suffice. Integer solutions of the independence and absorption
conditions are easy to describe; identifying their convex hull is the
additional requirement. A fractional point satisfying the displayed
constraints need not be an average of actual kernels. Merely finding
one kernel also does not describe this polytope.
\subsection{Short English statement}
Describe the convex hull of kernel incidence vectors by explicit graph-based inequalities for suitable classes of digraphs.
\subsection{Sources}
\begin{itemize}
\item[S1] ULAM.AI and the UnsolvedMath Contributors, supplied problems.json, record 3000066 (AMR-029-0066), AMR Open Problem Lists. Catalogue identity and source transcription; CC BY 4.0.
\url{https://huggingface.co/datasets/ulamai/UnsolvedMath}.
\item[S2] Egres Open - List of open problems, Open-problem page: Polyhedral description of kernels. Original source indexed by AMR-029-0066.
\url{https://oldlemon.cs.elte.hu/egres/open/List_of_open_problems}.
\item[S3] EGRES Open, Polyhedral description of kernels. Individual source statement and remarks.
\url{https://lemon.cs.elte.hu/egres/open/Polyhedral_description_of_kernels}.
\end{itemize}
\end{document}
